% Vorgaben Assignment aus Studienheft SQL03
% Formatvorgaben fuer den Text
% Umfang: 8 - 10 Seiten (inkl. Abbildungen und Tabellen, aber ohne Deckblatt, % Gliederung und Literaturverzeichnis, Eidesstattliche Erklaerung)
% Zeilenabstand: 1,5
% Schriftart: frei
% Schriftgrad: 12 pt
% Variablen, physikalische Groessen und Funktionszeichen werden kursiv gedruckt.
% Korrekturrand: links: 4,5 cm, rechts 2,0 cm, oben und unten jeweils 3,0 cm
% Deckblatt: (Adresse, AKAD-E-Mail-Adresse, Immatrikulationsnummer, Modul-
% bezeichnung, Thema, Datum, Felder für Korrektor)
% Gliederung (1 Seite)
% Literaturverzeichnis (3 - 5 Literaturquellen  z. B. Lehrbuecher, aktuelle Fachartikel recherchieren)
% Eidesstattliche Erklaerung (unterschrieben und fest eingebunden)
% Bearbeitungsdauer: 2 Monate


\documentclass[a4paper,12pt]{article}
\usepackage[ngerman]{babel}
\usepackage[nottoc]{tocbibind} % Anzeigen des Literaturverzeichnisses im TOC
\usepackage{epsfig}
\usepackage{times}
\usepackage{supertabular}
\usepackage{wrapfig}
\usepackage{multirow}
\usepackage[onehalfspacing]{setspace}
\usepackage{listings}
\usepackage{mathptmx}
\usepackage{geometry}
\usepackage{helvet}
\usepackage{courier}
\usepackage{setspace}
\usepackage{textcomp}
\usepackage[T1]{fontenc}
\usepackage[utf8]{inputenc}
\usepackage{fancyhdr}
\usepackage{float} % Notwendig fuer figure[h]
\usepackage[printonlyused]{acronym}

\newif\iflistoffigures
\newif\iflistoftables
\newif\ifacronym

%Titel
\newcommand*{\Titel}{Intranet der Dinge} 

%Betreff
\newcommand*{\Betreff}{Assignment} 

%Betreuer
\newcommand*{\Betreuer}{Prof. J. Anton Illig} 

%Vor- und Nachname
\newcommand*{\Name}{Stefan Waidele}

%Straße und Hausnummer
\newcommand*{\Strasse}{Ensisheimer Straße 2} 

%Plz und Ort
\newcommand*{\PlzOrt}{79395 Neuenburg am Rhein} 

%Immatrikulationsnummer
\newcommand*{\Immatrikulationsnummer}{102 81 71}

%Email 
\newcommand*{\Email}{Stefan@Waidele.info} 


% Verzeichnisse (Wenn nicht benötigt, Zeile mit % auskommentieren oder löschen

%% Abbildungsverzeichnis 
%\listoffigurestrue
%% Tabellenverzeichnis
%\listoftablestrue
%% Abkürzungsverzeichnis
%\acronymtrue

% Wittwen und Waisen verhindern
\clubpenalty10000
\widowpenalty10000
\displaywidowpenalty=10000

\usepackage[flushmargin,hang,ragged]{footmisc}
\usepackage{lmodern} %Type1-Schriftart für nicht-englische Texte
\usepackage{fancyhdr}

\usepackage[
	pdftitle={\Titel},
	pdfsubject={WIN03 -- Innovative Themen der Wirtschaftsinformatik},
	pdfauthor={Stefan Waidele},
	pdfkeywords={akad, win03, assignment, wirtschaftsinformatik}
	hyperfootnotes=false,
	colorlinks=true,
	linkcolor=black,
	urlcolor=black,
	citecolor=black
]{hyperref}

%\renewcommand{\familydefault}{\rmdefault}
\renewcommand{\bflabel}[1]{\normalfont{\normalsize{#1}}\hfill}

%% Definition for Codeschnipsel im Fließtext
\newcommand{\code}{\texttt}
\newcommand{\buzz}{\textit}
\newcommand{\todo}[1]{\fbox{\parbox{\textwidth}{\textbf{To do:} #1}}}
%\newcommand{\myref}[1]{„\ref{#1}~\nameref{#1}“}
\newcommand{\myref}[1]{\textit{\ref{#1}~\nameref{#1}}}

%% Für Codeblöcke mit Syntax-Highlighting
%% http://www.ctan.org/tex-archive/macros/latex/contrib/minted/
\usepackage{minted}
\definecolor{bg}{rgb}{0.95,0.95,0.95}


\makeatother

\geometry{a4paper, left=45mm, right=20mm, top=30mm, bottom=30mm}
\pagenumbering{roman}

\begin{document}
	
\parskip=1em
\parindent=0cm

%description: Deckblatt in Deutsch
%% Basierend auf einer TeXnicCenter-Vorlage von Tino Weinkauf.
%% sowie der akad-vorlage von Daniel Falkner
%%%%%%%%%%%%%%%%%%%%%%%%%%%%%%%%%%%%%%%%%%%%%%%%%%%%%%%%%%%%%%

%%%%%%%%%%%%%%%%%%%%%%%%%%%%%%%%%%%%%%%%%%%%%%%%%%%%%%%%%%%%%
%% Deckblatt
%%%%%%%%%%%%%%%%%%%%%%%%%%%%%%%%%%%%%%%%%%%%%%%%%%%%%%%%%%%%%
%%
%% ACHTUNG: Sie benötigen ein Hauptdokument, um diese Datei
%%          benutzen zu können. Verwenden Sie im Hauptdokument
%%          den Befehl "\input{dateiname}", um diese
%%          Datei einzubinden.
%%

\begin{titlepage}

%\thispagestyle{empty}

\vspace{5cm}

\Name \\ 
\Strasse \\ 
\PlzOrt\\ 
\href{mailto:\Email}{\Email}

AKAD University\\
Immatrikulationsnummer: \Immatrikulationsnummer

\vfill

\begin{tabbing}
Modul WIN03 --- \=Innovative Themen der Wirtschaftsinformatik\\ 
				\>Assignment
\end{tabbing}

\LARGE
\textsc{Web 3.0 \\Daten sind das Öl des 21. Jahrhunderts}\\

\vfill

\normalsize

Betreuer: \Betreuer

Entwurf vom \today %%Datum der Abgabe - am besten selbst reinschreiben.

\vfill

\includegraphics[width=3cm]{akad_logo.png}

AKAD University

\end{titlepage}

\normalsize


\normalsize

\begin{spacing}{1.0} % Verzeichnisse werden mit einzeiligem Abstand gesetzt
\parskip=0em
\newpage

% Inhaltsverzeichnis
\setcounter{tocdepth}{2}
\tableofcontents 
\newpage

% Abbildungsverzeichnis
\iflistoffigures
\listoffigures 
\newpage
\fi

% Tabellenverzeichnis
\iflistoftables
\listoftables 
\newpage
\fi

% Abkürzungsverzeichnis
\ifacronym
\section*{Abkürzungsverzeichnis}
\addcontentsline{toc}{section}{Abkürzungsverzeichnis} 
\begin{acronym}[\hspace{5cm}]
	\acro{CAPTCHA}{Completely Automated Public Tourint Test to Tell Computers and Humans Apart}
	\acro{CSS}{Cascading Style Sheets}
	\acro{DoS}{Denial of Service}
	\acro{DDoS}{Distributed Denial of Service}
	\acro{ERP}{Enterprise Resource Planning}
	\acro{HTML}{Hypertext Markup Language}
	\acro{HTTP}{Hypertext Tranfsfer Protocol}
	\acro{HTTPS}{Hypertext Tranfsfer Protocol Secure}
	\acro{OWASP}{Open Web Application Security Project}
	\acro{PDO}{PHP Data Objects}
	\acro{SCM}{Supply Chain Management}
	\acro{SQL}{Structured Query Language}
	\acro{WWW}{World Wide Web}
	\acro{XSS}{Cross Site Scripting}
%%
%	
%	\acro{}{}
%	\acro{}{}
%	\acro{}{}
%	\acro{}{}
%	\acro{}{}
%	\acro{}{}
%	\acro{}{}
%	\acro{}{}
	\acro{Test}{Wird nicht im Text verwendet und taucht auch nicht im Verzeichnis auf}
\end{acronym}

\fi

\parskip=1em
\end{spacing} 

\clearpage

\newcounter{romanPagenumber} 
\setcounter{romanPagenumber}{\value{page}} % Roemische Seitenanzahl speichern.


\pagestyle{fancy}
\fancyhead{}
\fancyhead[LO,RE]{\textsc{\Titel}}
\fancyhead[RO,LE]{\thepage}
\fancyfoot[CO,CE]{}
\setlength{\headheight}{15pt}

% \nocite{*} 

\pagenumbering{arabic}

\begin{spacing}{1.5} % Zeilenabstand: 1,5 fuer den Textteil

\section{Einleitung}
\label{sec:einleitung}

\subsection{Begründung der Problemstellung}

Durch Entwickungen wie \buzz{Internet of Things} und \buzz{Ubiquitous Computing} werden die in naher Zukunft die generierte Datenmenge als auch die Anzahl der verarbeitenden Instanzen in den nächsten Jahren wohl deutlich zunehmen.

\subsection{Ziele dieser Arbeit}

\textbf{Ziel dieser Arbeit ist es, die momentanen Entwicklungen der Datenbeständen hin zum semantischen Web und deren auswirkungen auf Wirtschaft und Gesellschaft mit den durch die Entdeckung des Öls im 20. Jahrhundert zu vergleichen.}

Hierzu werden zunächst im Kapitel~\myref{sec:grundlagen} die für diese Arbeit relevanten Begriffe und Konzepte definiert, bevor im Kapitel~\myref{sec:technologien} die momentan verfügbaren Technologien genannt und erklärt werden. 

Darauf aufbauend werden im Kapitel~\myref{sec:herausforderungen} die technischen Probleme identifiziert, und Lösungsansätze skizziert. Im Kapitel~\myref{sec:potentiale} werden schließlich einige mit dem IoT erwachsenden Möglichkeiten beschrieben. 

\todo{An die tatsächliche Arbeit anpassen}

\subsection{Abgrenzung}

\todo{An die tatsächliche Arbeit anpassen}

\section{Grundlagen}
\label{sec:grundlagen}

\subsection{Veredlungsprozess von Öl}

Rohöl: Finden / Fördern\\
Raffinerie: Sammeln / Aufbereiten\\
Anwendung: Verschiedene Ölprodukte, aber auch Kunststoffe, Kaugummi, etc.

\subsection{Auswirkungen des Öls auf die Wirtschaft des 20. Jahrhunderts}

Für Unternehmen und Verbraucher

\subsection{Definition: Daten, Informartion, Wissen}

\subsection{Semantisches Markup}

Beschreibung der Daten außerhalb der Datei

\subsection{Metadaten}

Beschreibung der Daten in der Datei (z.B. EXIF--Tags)

\subsection{Automatische Informationsgewinnung}

z.B. Gesichtserkennung

\subsection{Automatische Wissensgewinnung}

Durch Verknüpfung von Informationen kann Wissen generiert werden:\\
Person A ist auf einem Bild zusammen mit Person B zu sehen (Gesichtserkennung). Die Geoinformation und Uhrzeit (EXIF--Tags) zeigen, dass das Foto auf einer Veranstaltung aufgenommen wurde, das zu diesem Zeitpunkt an diesem Ort statt fand (Semantisches Markup der Veranstaltung).

Daraus kann auf politische Gesinnung oder Trinkfreudigkeit von Person A geschlossen werden.
\section{Auswirkungen des Semantischen Webs auf die Wirtschaft}

\subsection{}

\section{Auswirkungen des Semantischen Webs auf die Gesellschaft}

\subsection{}

\section{Vergleich der Auswirkungen mit denen des Öls}

\subsection{}

\section{Fazit \& Ausblick}

\subsection{Fazit}

Datenschutz vergleichbar mit Umweltschutz (Öl)

\subsection{Ausblick}

Die Entwicklungen des Web x.0 scheinen sich immer mehr von Big Data über Big Information bis hin zu Big Knowledge hin entwickeln. Diese Entwicklung hat durchaus das Potential, ähnliche oder sogar größere Auswirkungen auf Wirtschaft und Gesellschaft haben, wie das Erdöl im 20. Jahrhundert hatte.
\include{anhang}

\end{spacing}

\clearpage

% Literaturverzeichniss - Ab hier wieder Roemische Seitenzahlen
\pagestyle{plain}
\pagenumbering{roman}
\setcounter{page}{\theromanPagenumber}
%\addcontentsline{toc}{section}{Literatur-- und Quellenverzeichnis}
\renewcommand{\refname}{Literatur-- und Quellenverzeichnis}
\bibliographystyle{apalike}
\bibliography{literatur}
\onehalfspacing
\clearpage

\pagestyle{empty} 
\thispagestyle{empty}

\begin{center}
{\Large Eidesstattliche Erklärung}
\vspace*{4cm}\end{center}
\noindent
Ich versichere, dass ich das beiliegende Assignment selbstständig verfasst, keine anderen als die angegebenen Quellen und Hilfsmittel benutzt sowie alle wörtlich oder sinngemäß übernommenen Stellen in der Arbeit gekennzeichnet habe. 
\vspace{3cm}

\hspace{-0.8cm}
\rule[0.5ex]{6.5cm}{1pt}
\hspace{1.3cm}
\rule[0.5ex]{6.5cm}{1pt}
\\(Datum, Ort)
\hspace{6.3cm}
(Unterschrift)

\clearpage
\end{document}

